\subsection{局部凸空间上的 Mackey 拓扑与弱化拓扑}

设 E 为局部凸空间，\(E'\) 为其对偶。我们将 E 与 \(E'\) 置于对偶，所借助的典范双线性型为 \((x, x') \mapsto \langle x, x' \rangle\)，它定义在 \(E \times E'\) 上。这一对偶在 \(E'\) 中是分离的。我们将考虑 E 上与 E 和 \(E'\) 之间的对偶相容的三种拓扑：

\begin{enumerate}
\def\labelenumi{\alph{enumi})}
\item
  E 上原先给定的拓扑，当可能产生混淆时，我们称其为初始拓扑；
\item
  拓扑 \(\sigma(\mathrm{E},\mathrm{E}^{\prime})\)，称为 \(\mathrm{E}\) 上的弱化拓扑；
\item
  拓扑 \(\tau(\mathrm{E},\mathrm{E}')\)，称为 E 上的 Mackey 拓扑。
\end{enumerate}

初始拓扑比弱化拓扑细，而比 Mackey 拓扑粗；此外，这三种拓扑可以互不相同（IV，第 49 页，习题 8）。

由 IV 第 1 页的命题 1，这三种拓扑具有相同的闭凸集、相同的桶、相同的有界集以及相同的适配有界结构。特别地：

命题 2. --- 设 E 为局部凸空间，A 为 E 的一个凸子集（例如 E 的一个向量子空间）。A 对于 E 的初始拓扑的闭包与对于 E 的弱化拓扑的闭包相同。

注. --- 1) E 中元素的一族 \((x_{i})_{i\in\mathbf{I}}\) 对于初始拓扑是完全的（相应地，拓扑无关的），当且仅当它对于弱化拓扑也是如此；这由命题 2 得出。因此我们可以应用 II 第 43 页的判别法。

\begin{enumerate}
\def\labelenumi{\arabic{enumi})}
\setcounter{enumi}{1}
\tightlist
\item
  设 \(\mathcal{T}_1\) 与 \(\mathcal{T}_2\) 是 E 上的两个局部凸拓扑，与 E 和 E' 之间的对偶相容，且 \(\mathcal{T}_1\) 比 \(\mathcal{T}_2\) 细。则 0 的每个对于 \(\mathcal{T}_1\) 的邻域，只要它是凸的且对于 \(\mathcal{T}_1\) 闭，就对于 \(\mathcal{T}_2\) 闭（依据 IV 第 1 页的命题 1）。因此（GT, II, § 3, 第 3 目，推论）E 中每个对于 \(\mathcal{T}_2\) 完备的子集对于 \(\mathcal{T}_1\) 也完备。
\end{enumerate}

特别地，E 中每个对于弱化拓扑完备的子集对于初始拓扑完备，而 E 中每个对于初始拓扑完备的子集对于 Mackey 拓扑也完备。若 E 对于弱化拓扑拟完备，则它对于每个与 E 和 E' 之间的对偶相容的拓扑都如此。若它对于初始拓扑拟完备，则它对于 Mackey 拓扑也如此。

\begin{enumerate}
\def\labelenumi{\arabic{enumi})}
\setcounter{enumi}{2}
\tightlist
\item
  假定 E 是 Hausdorff 的（对初始拓扑而言）。设 A 是 E 的一个子集，它对于 \(\sigma(\mathrm{E}, \mathrm{E}')\) 闭且有界，从而对于每个与 E 和 \(E'\) 之间的对偶相容的拓扑也如此。由于 A 对于 \(\sigma(\mathrm{E}, \mathrm{E}')\) 预紧（III，第 3 页，注 5），假定 A 对于 \(\sigma(\mathrm{E}, \mathrm{E}')\) 紧等价于 A 对于 \(\sigma(\mathrm{E}, \mathrm{E}')\) 完备。
\end{enumerate}

因此，根据注 2，我们看出：

命题 3. --- 假定 E 是 Hausdorff 的，E' 是它的对偶。E 中每个对于初始拓扑预紧且对于 \(\sigma(E, E')\) 紧的子集，对于初始拓扑是紧的。

\begin{enumerate}
\def\labelenumi{\arabic{enumi})}
\setcounter{enumi}{3}
\tightlist
\item
  拓扑 \(\beta(\mathrm{E},\mathrm{E}')\)（IV，第 4 页，定义 3）比 Mackey 拓扑细。若 \(\beta(\mathrm{E},\mathrm{E}')\) 不同于 \(\tau(\mathrm{E},\mathrm{E}')\)，则它与 E 和 E' 之间的对偶不相容。E 是桶型空间当且仅当初始拓扑等于 \(\beta(\mathrm{E},\mathrm{E}')\)（III，第 24 页）。
\end{enumerate}

命题 4. --- 设 E 为局部凸空间。在下列每种情形中，E 上的 Mackey 拓扑都与初始拓扑相同：

\begin{enumerate}
\def\labelenumi{\alph{enumi})}
\item
  E 是桶型空间；
\item
  E 是有界型空间；
\item
  E 是可度量化的。
\end{enumerate}

我们首先注意到，E 上的 Mackey 拓扑与初始拓扑相同，当且仅当 \(E'\) 的每个对于 \(\sigma(E', E)\) 紧的凸子集都是等度连续的。当 E 是桶型空间时情形当然如此（III，第 24 页，推论）。

设 E 是有界型空间；令 V 是 E 中对于拓扑 \(\tau(E, E')\) 的一个凸均衡零邻域。令 B 是 E 的一个对于初始拓扑有界的子集。由于 B 对于 Mackey 拓扑有界，V 吸收 B；又由于 E 是有界型空间，V 是对于初始拓扑的零邻域。

在情形 c)，空间 E 是有界型的（III，第 12 页，命题 2）。

